% Compilar dos veces: pdflatex investigacion_robotica.tex
\documentclass[11pt,a4paper]{article}
\usepackage[utf8]{inputenc}
\usepackage[T1]{fontenc}
\usepackage{lmodern}
\IfFileExists{spanish.ldf}{\usepackage[spanish,es-tabla]{babel}}{}
\usepackage[margin=2.4cm]{geometry}
\usepackage{microtype}
\usepackage{booktabs,longtable,array}
\usepackage{xcolor}
\usepackage{tabularx}
\usepackage{graphicx,float,caption}
\graphicspath{{imagenes/}{output/imagenes/}}
\captionsetup{font=small,labelfont=bf,justification=raggedright,singlelinecheck=false}
\newcommand{\robotfoto}[4]{%
  \begin{figure}[H]
  \centering
  \includegraphics[width=#2,height=8.8cm,keepaspectratio]{#1.jpg}
  \caption{#3}\label{fig:#1}
  {\footnotesize\raggedright #4.\par}
  \end{figure}}

\usepackage{xurl}
\usepackage{hyperref}
\definecolor{azul}{HTML}{173E60}
\hypersetup{colorlinks=true,linkcolor=azul,citecolor=azul,urlcolor=azul,
  pdftitle={Aplicaciones de los robots: asistencia, servicio y robótica social},
  pdfsubject={Investigación documental con casos empresariales y condiciones sociales}}
\AtBeginDocument{\renewcommand{\abstractname}{Resumen}\renewcommand{\refname}{Referencias}\renewcommand{\figurename}{Figura}\renewcommand{\tablename}{Cuadro}}
\setlength{\parindent}{0pt}
\setlength{\parskip}{0.55em}
\setlength{\emergencystretch}{2em}
\renewcommand{\arraystretch}{1.18}
\newcolumntype{P}[1]{>{\raggedright\arraybackslash}p{#1}}
\title{Aplicaciones de los Robots}
\author{Dan Alvarado Solano \and Gabriel Molina Bulgarelli}
\date{}

\begin{document}
\maketitle
\begin{abstract}
TODO
\end{abstract}

\section{Introducción}
Los robots se aplican a tareas y entornos diversos. Este documento estudia sus usos de asistencia, de servicio y de interacción social, atendiendo al propósito de cada sistema y a las personas que participan en su funcionamiento.

\section{Alcance y fuentes de la investigación}
Se examinan ejemplos de compañías y productos reconocibles, junto con las condiciones que justifican estudiar su uso. La clasificación se organiza según la función principal y el contexto de aplicación.

\section{Desarrollo: características de los robots de servicio}
TODO

\subsection{Precisiones sobre autonomía y clasificación}
TODO

\section{Aplicaciones generales y empresas representativas}
TODO

\subsection{ABB Robotics y la manufactura}
TODO

\subsection{Spot y la inspección de instalaciones}
TODO

\section{Robótica de asistencia: autonomía cotidiana y rehabilitación}
TODO

\subsection{Jaco, de Kinova}
Jaco es un brazo robótico que se instala en una silla de ruedas motorizada para personas con movilidad limitada o ausente en las extremidades superiores. Kinova describe tareas como comer, beber, recoger objetos, abrir puertas y accionar interruptores. El usuario puede controlarlo mediante el mando de la silla u otras interfaces adaptadas, incluidos controles de cabeza o de soplo y aspiración.


\textbf{Cómo se desarrolló.} Un trabajo técnico de investigadores de Kinova y de la Universidad Laval describe la evolución del brazo y de su control. La estructura combina seis segmentos de fibra de carbono con actuadores rotativos. El desarrollo se centró también en facilitar acciones: posiciones de alcance y plegado, filtros para suavizar el movimiento y coordinación de traslación y orientación al beber. Esas decisiones reducen la cantidad de movimientos que el usuario debe combinar mediante la interfaz de su silla.






Para evaluar su utilidad, hay que observar qué actividad puede completar la persona, cuánto esfuerzo requiere y cuándo sigue necesitando ayuda. Poder recoger un teléfono o acercar un vaso puede modificar una rutina cotidiana. La prueba del equipo debe realizarse con la persona y en su entorno. La descripción comercial acredita las funciones previstas, pero no demuestra que todos los usuarios obtengan el mismo resultado.

\subsection{HAL, de CYBERDYNE}
HAL es una familia de dispositivos de asistencia que se llevan sobre el cuerpo. En la versión médica para miembros inferiores, sensores sobre la piel detectan señales bioeléctricas relacionadas con el intento de movimiento y el sistema ayuda a mover las articulaciones. CYBERDYNE lo presenta para tratamiento de alteraciones de la marcha y advierte que las indicaciones autorizadas varían según el país y la autoridad reguladora.


\textbf{Cómo se desarrolló.} Yoshiyuki Sankai inició esta línea de investigación en la Universidad de Tsukuba, combinando robótica, señales del cuerpo y control del movimiento. La reseña de la Organización Mundial de la Propiedad Intelectual documenta la creación de CYBERDYNE en 2004 para desarrollar y comercializar los resultados. La transferencia requirió prototipos, propiedad intelectual y fabricación, además del mecanismo de asistencia. El propósito técnico fue acoplar la ayuda del dispositivo al intento de movimiento de la persona, en lugar de imponer únicamente una secuencia fija.






En este caso conviene separar dos preguntas: si el dispositivo ayuda a ejecutar el movimiento durante una sesión y si el tratamiento produce una mejoría funcional que continúa sin el equipo. Son resultados distintos. La información del fabricante permite explicar el mecanismo, pero no basta para afirmar recuperación permanente ni para recomendarlo a cualquier persona con dificultad para caminar.

\section{Robótica de servicio: tareas y despliegues públicos}
TODO

\subsection{Starship y el reparto en espacios peatonales}
Starship Technologies desarrolla robots para entregar comida y compras. Su documentación explica que utilizan rutas peatonales y que disponen de asistencia humana remota cuando la necesitan. La autonomía de navegación, por tanto, convive con una estructura de apoyo humano.


\textbf{Cómo se desarrolló.} La cronología de la compañía sitúa su fundación por Ahti Heinla y Janus Friis en julio de 2014 y el primer prototipo en agosto de ese año. Después vinieron pruebas en ciudades europeas, un piloto de reparto de comida en Londres en 2017 y el servicio comercial de compras en Milton Keynes en 2018. El paso de prototipo a servicio requirió probar rutas y encuentros con peatones, además de organizar la entrega y el apoyo remoto ante excepciones.






En un comunicado del 28 de abril de 2026, la empresa anunció diez millones de entregas y una flota superior a tres mil robots en ocho países. Son cifras declaradas por el proveedor, útiles para documentar la escala anunciada del servicio. No constituyen una auditoría independiente de seguridad o rentabilidad. La evaluación social debe incluir a quienes comparten la acera: se propone observar bloqueos, incidentes, accesibilidad y rapidez de respuesta ante un equipo detenido.

\subsection{Moxi y la logística hospitalaria}
Moxi, de Diligent Robotics, realiza entregas dentro del hospital. La documentación de la empresa describe compartimentos con cierre para muestras, medicamentos y suministros, así como el uso de ascensores y puertas. También indica que el robot puede recibir solicitudes de entrega del personal y volver a su estación de carga.


\textbf{Cómo se desarrolló.} Andrea Thomaz y Vivian Chu fundaron Diligent Robotics en 2017. La empresa identifica como bases de su trabajo la manipulación móvil, la interacción con personas y el aprendizaje guiado por humanos. Moxi combina una base que se desplaza con un brazo para manipular objetos; su uso hospitalario exige integrar solicitudes, puertas, ascensores y puntos de entrega. Así, el desarrollo incluye tanto el robot como su conexión con las rutinas del centro.






El objetivo operativo es reducir desplazamientos que consumen tiempo del personal. Sin embargo, contar entregas no permite concluir que mejora la atención clínica. Una evaluación local debería medir demoras, errores, entregas completadas y tiempo neto disponible, incluyendo el trabajo que exige atender fallos del robot. Estas son propuestas de medición del informe, no resultados demostrados por la página comercial.


El informe mundial de enfermería de 2025, difundido por la OMS, estimó un déficit de 5,8 millones de profesionales en 2023 y describió fuertes desigualdades entre países. Ese déficit ayuda a entender el interés por automatizar tareas de transporte, pero no prueba que un robot resuelva la falta de personal. La propia OMS señala la necesidad de invertir en educación, empleo y retención.

\subsection{UVD Robots durante la pandemia de COVID-19}
La pandemia ofrece un caso histórico de despliegue público. En febrero de 2021, la Comisión Europea informó sobre la llegada de robots de desinfección de la compañía danesa UVD Robots a hospitales europeos. Los equipos utilizaban luz ultravioleta y el personal iniciaba la operación desde fuera de la habitación, sin personas presentes.


\textbf{Cómo se desarrolló.} El desarrollo antecede a la pandemia: la Comisión Europea documenta un proyecto de 2014 a 2017 impulsado mediante contratación precomercial de las regiones danesas. Blue Ocean Robotics trabajó en una plataforma móvil de desinfección; UVD Robots se constituyó en 2016 dentro de esa colaboración entre sector público y empresa. La solución combinó navegación autónoma y lámparas UV-C. El caso muestra cómo una necesidad hospitalaria puede orientar prototipos y pruebas antes de una compra a mayor escala.






En diciembre de 2021, la Comisión comunicó que había completado la entrega de 305 robots de desinfección a hospitales de los 27 Estados miembros. Esta documentación demuestra una respuesta institucional y un despliegue durante una emergencia sanitaria. Para afirmar cuánto redujeron las infecciones se necesitarían resultados clínicos comparativos; la cantidad de equipos entregados no responde esa pregunta.

\subsection{Da Vinci y el papel del equipo humano}
TODO

\section{Robótica social: comunicación, acompañamiento y evidencia}
TODO

\subsection{Pepper y la atención al público}
Pepper, desarrollado por SoftBank y Aldebaran, utiliza conversación, movimiento y una pantalla para interactuar con personas. La documentación de Aldebaran describe usos de recepción, información y orientación. Se presenta como un caso histórico documentado de interacción pública, cuya incorporación requiere verificar la disponibilidad y el soporte local.


\textbf{Cómo se desarrolló.} SoftBank y Aldebaran presentaron Pepper en junio de 2014, después de que SoftBank invirtiera en Aldebaran en 2012. La plataforma reunió motores para gestos, cámaras, micrófonos y una pantalla, junto con herramientas de desarrollo basadas en NAOqi. El diseño permitió crear aplicaciones que combinan habla, movimientos y contenido visual. La presentación comercial describía inferencias a partir del rostro y la voz; esas funciones deben entenderse como procesamiento de señales, sin atribuir al robot sentimientos o comprensión humana.






La apariencia y la capacidad de responder pueden atraer la atención de visitantes, pero el resultado relevante es si la persona consigue la información o la ayuda que busca. Para una recepción se propone medir consultas resueltas, errores y derivaciones al personal. Las expresiones o respuestas programadas del robot no bastan para atribuirle comprensión emocional humana.

\subsection{ElliQ y la soledad en personas mayores}
ElliQ, de Intuition Robotics, es un dispositivo de acompañamiento que inicia conversaciones y propone actividades. Un informe de febrero de 2026 alojado por la Oficina para el Envejecimiento del Estado de Nueva York (NYSOFA), elaborado por Intuition Robotics, describe el programa conjunto y presenta datos del período de junio de 2024 a mayo de 2025. Según ese informe, el 94\,\% de los clientes manifestó sentirse menos solo.


\textbf{Cómo se desarrolló.} Intuition Robotics anunció la disponibilidad comercial en marzo de 2022, tras más de cinco años de investigación, desarrollo y pruebas con usuarios, según su comunicado. El diseño realizado con fuseproject separó una unidad que expresa orientación mediante movimiento y luz de una pantalla para contenido visual. Su forma abstracta y la iniciativa para conversar son decisiones de interacción: el proyecto buscó una presencia doméstica que propusiera actividades y acompañara rutinas. Las pruebas y revisiones del producto no equivalen, por sí mismas, a una validación clínica.






El porcentaje corresponde a una valoración declarada por los clientes; no significa que la soledad haya disminuido un 94\,\%. El documento tampoco presenta ese dato como resultado de un ensayo aleatorizado con grupo de control. Por ello, este informe lo interpreta como información del programa sobre la experiencia de uso, sin atribuirle una demostración causal de eficacia clínica. La participación del proveedor en la elaboración exige mantener explícito el origen de los resultados.

\subsection{PARO y el cuidado de personas con demencia}
PARO es un robot con forma de cría de foca desarrollado en el instituto japonés AIST y comercializado por Intelligent System. Su interacción permite incorporarlo a actividades de cuidado y acompañamiento. Su uso en cuidados se ha estudiado mediante ensayos comparativos publicados.


\textbf{Cómo se desarrolló.} AIST documenta que Takanori Shibata comenzó esta investigación en 1993. En 2004 presentó un prototipo de octava generación, desarrollado en colaboración con Microjennics, y una licencia comercial para Intelligent System. Las iteraciones integraron sensores de contacto, luz, sonido, temperatura y postura para producir respuestas durante la interacción. La elección de una cría de foca dio lugar a una forma que se puede sostener y acariciar; el desarrollo técnico se centró en esas respuestas y en la adaptación al contacto.






Moyle y colaboradores estudiaron a 415 residentes con demencia en 28 centros de cuidados de larga duración de Australia. El ensayo aleatorizado por centros comparó PARO, un peluche de aspecto similar y la atención habitual. PARO favoreció la participación verbal y visual frente al peluche. Aunque algunas observaciones favorecieron a PARO frente a la atención habitual en estado de ánimo y agitación, la escala CMAI-SF no mostró diferencias entre grupos en agitación. La superioridad frente al peluche se limitó al aumento de la participación.


Estos resultados respaldan una aplicación delimitada y muestran por qué importa elegir el comparador. Una mayor participación durante la actividad no demuestra mejoría general de la demencia. Tampoco permite trasladar automáticamente el resultado a otros robots o poblaciones.

\section{Comparación entre categorías}
TODO

\section{Condiciones sociales que hacen visibles estas aplicaciones}
TODO

\subsection{Discapacidad y acceso}
TODO

\subsection{Envejecimiento y necesidades de apoyo}
TODO

\subsection{Soledad y aislamiento social}
TODO

\subsection{Desigualdad, trabajo y emergencias}
TODO

\section{Criterios para una incorporación responsable}
TODO

\section{Conclusiones}
TODO

\begin{verbatim}
https://www.softbank.jp/en/corp/group/sbm/news/press/2014/20140605_01/
https://aging.ny.gov/system/files/documents/2026/02/nysofa-elliq-project-update-2026.pdf
https://doi.org/10.1109/ICORR.2005.1501143
https://aldebaran.com/es/pepper
https://ifr.org/service-robots
https://www.diligentrobots.com/about/
https://digital-strategy.ec.europa.eu/en/news/305-disinfection-robots-delivered-hospitals-across-eu
https://digital-strategy.ec.europa.eu/en/news/coronavirus-first-eu-disinfection-robots-arrive-hospitals
https://uvd.blue-ocean-robotics.com/about
https://www.who.int/news/item/12-05-2025-nursing-workforce-grows--but-inequities-threaten-global-health-goals
https://assistive.kinovarobotics.com/product/jaco-robotic-arm
https://www.starship.xyz/about/
https://www.idsa.org/awards-recognition/idea/idea-gallery/elliq-the-sidekick-for-healthier-happier-aging/
https://doi.org/10.1016/S1697-7912(08)70140-7
https://www.diligentrobots.com/faq
https://www.intuitionrobotics.com/post/announcement-commercial-availability-of-elliq
https://www.starship.xyz/press/autonomous-delivery-moves-into-the-mainstream-as-starship-technologies-passes-10-million-deliveries/
https://www.resna.org/sites/default/files/conference/2016/other/campeau_lecours.html
https://doi.org/10.1016/j.jamda.2017.03.018
https://www.aist.go.jp/aist_e/list/latest_research/2004/20041208_2/20041208_2.html
https://digital-strategy.ec.europa.eu/en/news/european-commission-procurement-triggers-wider-deployment-innovative-disinfection-robots-across
https://www.cyberdyne.jp/en/products/hal-lowerlimb-medical
https://www.starship.xyz/faq/
https://www.wipo.int/en/web/ip-advantage/w/stories/rise-of-the-cybernoids-japanese-inventor-creates-wearable-robots
\end{verbatim}


\end{document}
